% ตัวอย่างเฉลยที่ถูกต้อง 1 ข้อ — ใช้เป็นแม่แบบว่า "ครบ 5 ส่วน" หน้าตาเป็นอย่างไร
% ห้าม \input ไฟล์นี้เข้าไปในเฉลยจริง เปิดอ่านอย่างเดียว
%
% โจทย์: เลือกกรรมการ 3 คนจากสมาชิก 8 คน ได้กี่วิธี
% ตัวเลือกทั้งสี่มาจาก verify.py:
%   R[1] = C(8,3) = 56
%   D[1] = {'คูณ n กับ r ตรง ๆ': 24, 'หารด้วย r แทน r!': 112,
%           'นับลำดับด้วย (ใช้ P แทน C)': 336}
% ส่วน "ทำไมตัวเลือกอื่นผิด" ข้างล่างใช้ key ของ D[] ตรง ๆ ไม่ได้เขียนคำอธิบายใหม่

\Needspace{4\baselineskip}
\item
\ansline{ค.\ $56$}{}

\textbf{แนวคิด}~~ เลือกคน 3 คนจาก 8 คน โดยไม่สนลำดับที่เลือก จึงเป็นการจัดหมู่
(combination) ไม่ใช่การเรียงสับเปลี่ยน

\textbf{วิธีทำ}
\stepa{1}{สังเกตว่า "คณะกรรมการ" ไม่มีตำแหน่ง สลับชื่อกันแล้วได้คณะเดิม จึงไม่นับลำดับ}
\stepa{2}{สูตรจัดหมู่ $\displaystyle\binom{n}{r}=\frac{n!}{r!\,(n-r)!}$ ในที่นี้ $n=8$, $r=3$}
\stepa{3}{แทนค่า $\displaystyle\binom{8}{3}=\frac{8!}{3!\,5!}
          =\frac{8\times7\times6}{3\times2\times1}=\frac{336}{6}=56$}
\textbf{ตอบ} $56$ วิธี

\textbf{ทำไมตัวเลือกอื่นผิด}~~ ก. $24$ คือคูณ $n$ กับ $r$ ตรง ๆ
ข. $112$ คือหารด้วย $r$ แทน $r!$
ง. $336$ คือ $P(8,3)$ ของคนที่นับลำดับด้วย

\textbf{จุดพลาดที่พบบ่อย}~~ ใช้ $P(8,3)=336$ เพราะเผลอนับลำดับ ให้ถามตัวเองเสมอว่า
"สลับตำแหน่งกันแล้วเป็นคำตอบเดิมไหม" ถ้าเดิม = จัดหมู่

% สิ่งที่ตัวอย่างนี้ทำถูกและต้องทำตามทุกข้อ
%
% 1. เลือกวิธีที่ทำได้จริงในห้องสอบ — สูตรจัดหมู่ ไม่ใช่การไล่นับ 56 กรณี
% 2. เขียนสูตรรูปทั่วไปก่อนแทนค่า ไม่ใช่โผล่ 56 มาเฉย ๆ
% 3. ไม่ข้ามพีชคณิต — แสดง 336/6 ตรงกลาง ไม่กระโดดจากสูตรไปคำตอบ
% 4. ทุกตัวลวงมีที่มาจาก D[] และคำอธิบายคือชื่อความผิดที่คำนวณไว้จริง
%    ไม่ใช่ "ข. ผิดเพราะคำนวณผิด"
% 5. จบใน 3 ขั้นเพราะวิธีนี้ใช้ 3 ขั้นจริง ไม่ยืดให้ครบโควตา
%
% ตัวอย่างที่ผิดสำหรับข้อเดียวกัน — กระโดดขั้นแบบนี้คือสิ่งที่ต้องแก้:
%     \ansline{ค.\ $56$}{} จาก $\binom{8}{3}=56$
